\documentclass[11pt]{article}

% ── Packages ────────────────────────────────────────────────────────────────
\usepackage[margin=1in]{geometry}
\usepackage{amsmath, amssymb, amsthm}
\usepackage{mathtools}
\usepackage{enumitem}
\usepackage{hyperref}
\usepackage{parskip}

% ── Theorem environments ─────────────────────────────────────────────────────
\newtheorem{theorem}{Theorem}[section]
\newtheorem{lemma}[theorem]{Lemma}
\newtheorem{corollary}[theorem]{Corollary}
\newtheorem{proposition}[theorem]{Proposition}
\theoremstyle{definition}
\newtheorem{definition}[theorem]{Definition}
\newtheorem{example}[theorem]{Example}
\theoremstyle{remark}
\newtheorem{remark}[theorem]{Remark}

% ── Custom commands ───────────────────────────────────────────────────────────
\newcommand{\R}{\mathbb{R}}
\newcommand{\E}{\mathbb{E}}
\newcommand{\norm}[1]{\left\lVert#1\right\rVert}
\newcommand{\inner}[2]{\langle #1,\, #2 \rangle}
\newcommand{\T}{^{\top}}
\newcommand{\K}{K}  % number of classes
\DeclareMathOperator{\softmax}{softmax}

% ── Document ─────────────────────────────────────────────────────────────────
\begin{document}

\title{\textbf{Multiclass: Learning}\\[0.4em]
       \large Notes Template}
\author{}
\date{}
\maketitle
\tableofcontents
\newpage

% ─────────────────────────────────────────────────────────────────────────────
\section{Problem}
% ─────────────────────────────────────────────────────────────────────────────

\subsection{Setup}

% TODO: Describe the multiclass learning problem here.
% Example prompts:
%   - What is the model family (softmax regression, neural net, …)?
%   - What training data / likelihood model is assumed?
%   - What optimisation criterion is used (MLE, cross-entropy minimisation, …)?

\begin{definition}[Multiclass Learning Problem]
    % TODO: Fill in the formal definition.
    Given a training dataset
    $\mathcal{D} = \{(x^{(i)}, y^{(i)})\}_{i=1}^{n}$
    with $x^{(i)} \in \R^{d}$ and $y^{(i)} \in \{1, \ldots, \K\}$,
    the learning problem is to find parameters $\theta$ such that \ldots
\end{definition}

\subsection{Assumptions}

\begin{enumerate}[label=\textbf{A\arabic*.}]
    \item % TODO: State assumption 1 (e.g., parametric model for $p(y \mid x;\theta)$).
    \item % TODO: State assumption 2 (e.g., i.i.d.\ samples).
    \item % TODO: Add further assumptions as needed.
\end{enumerate}

\subsection{Objective}

% TODO: State the learning objective (e.g., maximum likelihood / cross-entropy loss).
\[
    % TODO: Replace with the actual objective.
    \hat{\theta} \;=\; \arg\min_{\theta} \;
        -\sum_{i=1}^{n} \log p\!\left(y^{(i)} \mid x^{(i)};\, \theta\right)
\]

% ─────────────────────────────────────────────────────────────────────────────
\section{Derivation}
% ─────────────────────────────────────────────────────────────────────────────

The following derivation establishes the optimal multiclass learning algorithm
step by step.

% ── Step 1 ────────────────────────────────────────────────────────────────────
\subsection{Step 1: [Title]}

% TODO: Describe what this step accomplishes (e.g., write out the joint likelihood).

\begin{proof}
    % TODO: Fill in the proof for Step 1.
    \[
        \text{(derivation goes here)}
    \]
\end{proof}

% ── Step 2 ────────────────────────────────────────────────────────────────────
\subsection{Step 2: [Title]}

% TODO: Describe what this step accomplishes (e.g., take the log and simplify).

\begin{proof}
    % TODO: Fill in the proof for Step 2.
    \begin{align}
        \ell(\theta) &= \sum_{i=1}^{n} \log p\!\left(y^{(i)} \mid x^{(i)};\,\theta\right) \notag \\
                     &= \text{(expand using model definition)} \label{eq:step2}
    \end{align}
\end{proof}

% ── Step 3 ────────────────────────────────────────────────────────────────────
\subsection{Step 3: [Title]}

% TODO: Describe what this step accomplishes (e.g., compute the gradient w.r.t.\ $\theta$).

\begin{proof}
    % TODO: Fill in the proof for Step 3.
    \[
        \nabla_{\theta}\, \ell(\theta) \;=\; \text{(gradient expression)}
    \]
\end{proof}

% ── Step 4 ────────────────────────────────────────────────────────────────────
\subsection{Step 4: [Title]}

% TODO: Describe what this step accomplishes
%       (e.g., derive the gradient update rule / show convexity).

\begin{proof}
    % TODO: Fill in the proof for Step 4.
    \[
        \theta \;\leftarrow\; \theta - \eta\, \nabla_{\theta}\, \ell(\theta)
        \quad \text{(placeholder update rule)}
    \]
\end{proof}

% ── Final Result ──────────────────────────────────────────────────────────────
\subsection{Final Result}

\begin{theorem}[Optimal Multiclass Learner]
    % TODO: State the final theorem.
    Under assumptions \textbf{A1}--\textbf{A3},
    the maximum-likelihood parameter estimate satisfies
    \[
        % TODO: Replace with the actual result (e.g., gradient condition, closed form).
        \nabla_{\theta}\, \ell\!\left(\hat{\theta}^{*}\right) \;=\; \mathbf{0}.
    \]
\end{theorem}

\begin{proof}
    Follows directly from Steps 1--4. \qed
\end{proof}

% ─────────────────────────────────────────────────────────────────────────────
\end{document}
